\chapter{Combinational Circuits}
\label{ch:comb}

In Chapter~\ref{ch:gates} we learned to write logic functions and build them
from gates, and in Chapter~\ref{ch:verilog} to describe a circuit as Verilog
text, but every time the function was handed to us ready-made. In
practice the task sounds different: ``sound the alarm if the door is open and
it is night or the guard is on.'' How do you turn such a sentence into a
circuit without using extra gates? And there is a second question left over
from the previous chapter: we treated gates as instantaneous, but a real
circuit needs time to ``settle down.'' What happens to the output during that
moment? This chapter answers both questions.

\begin{buildbox}
Open the \emph{Learn} tab, the book ``Logic gates,'' the lesson
\emph{``Two out of three''}: build a circuit in which lamp $y$ lights up if at
least two of the three buttons $a$, $b$, $c$ are pressed. Without reading
further, write down in your notebook how many blocks you used. At the end of
the chapter we will get the same circuit by the plan ``table $\to$
expression $\to$ simplification $\to$ circuit'' and compare.
\end{buildbox}

\section{Combinational and Sequential Logic}

All digital circuits fall into two large classes.

\begin{description}
  \item[Combinational circuits] --- the outputs depend \emph{only on the
    current values of the inputs}. They have no memory: apply the same inputs
    and you get the same outputs, no matter what happened before. Examples:
    logic gates, decoders, multiplexers, adders.
  \item[Sequential circuits] --- the outputs depend not only on the current
    inputs but also on the \emph{history}, that is, on what happened earlier.
    They have memory. Examples: flip-flops, registers, counters. We will get
    to them later; for now it is enough to know that they exist and that a
    combinational circuit cannot ``remember what happened.''
\end{description}

A feature that makes them easy to tell apart in a diagram: in a combinational
circuit the signal flows strictly from the inputs to the outputs, while in a
sequential circuit some of the outputs go back to the input through memory.

\begin{figure}[H]
\centering
\begin{tikzpicture}
  \node[block, minimum width=30mm, minimum height=20mm] (c) {Combinational\\logic};
  \foreach \i in {1,2,3} \draw[arr] ($(c.west)+(-1.2,0.9-\i*0.45)$) -- ($(c.west)+(0,0.9-\i*0.45)$);
  \foreach \i in {1,2} \draw[arr] ($(c.east)+(0,0.68-\i*0.45)$) -- ($(c.east)+(1.2,0.68-\i*0.45)$);
  \node[lbl, left=12mm of c] {inputs};
  \node[lbl, right=12mm of c] {outputs};
  \node[font=\sffamily\small\bfseries, above=2mm of c] {Combinational circuit};
  \begin{scope}[xshift=8.3cm]
    \node[block, minimum width=30mm, minimum height=20mm] (s) {Combinational\\logic};
    \node[gblock, minimum width=30mm, below=6mm of s] (m) {Memory\\(flip-flops)};
    \foreach \i in {1,2,3} \draw[arr] ($(s.west)+(-1.2,0.9-\i*0.45)$) -- ($(s.west)+(0,0.9-\i*0.45)$);
    \foreach \i in {1,2} \draw[arr] ($(s.east)+(0,0.68-\i*0.45)$) -- ($(s.east)+(1.2,0.68-\i*0.45)$);
    \draw[arr] ($(s.east)+(0.4,-0.7)$) -- ++(0,-1.0) |- (m.east);
    \draw[arr] (m.west) -- ++(-0.4,0) |- ($(s.west)+(0,-0.7)$);
    \node[lbl, left=4mm of m, align=center] {state};
    \node[font=\sffamily\small\bfseries, above=2mm of s] {Sequential circuit};
  \end{scope}
\end{tikzpicture}
\caption{The main difference: a sequential circuit has memory and feedback
through it}
\label{fig:comb_vs_seq}
\end{figure}

\begin{diybox}[Does the circuit have memory?]
\textbf{Task.} \stars{1} For each device decide whether it is combinational or
not: a lamp that ``is on while the button is held''; a lamp that ``turns on
with the first press and off with the second''; a ``sum of two numbers''
signal. \stars{2} Come up with an everyday device of each class.

\textbf{In Logra.} \emph{Designer} tab: build ``button $\to$ NOT $\to$
lamp'' and set the button to 0, then 1, then 0 again. Next open the example
\code{edge-toggle} (the example library, \emph{Code} tab) and hold the button
for several clock cycles.

\textbf{Expected result.} In the first circuit the same inputs always give
the same outputs. In the second the switching happens once per press: the
circuit ``remembers'' that the button has already been pressed.

\textbf{How to check.} \emph{Waveforms} tab: for a combinational circuit you
can predict the output from the truth table by looking at the current inputs;
for \code{edge-toggle} you cannot---you need the history as well.
\end{diybox}

\section{From a Word Description to a Circuit}

Designing a combinational circuit always follows the same plan:
\begin{enumerate}
  \item understand which inputs and outputs are needed and give them names;
  \item build the \textbf{truth table};
  \item write the logic expression;
  \item \textbf{simplify} it;
  \item draw the circuit.
\end{enumerate}

The main difficulty is the first step. A sentence in everyday language almost
never says how many inputs the problem has or what to do in ``borderline''
cases. So before the table, write down for each input what 1 means and what 0
means (``door open: 1,'' ``night: 1''), and only then build the table.

\subsection{Example: Three Judges Vote}

Three judges, $a$, $b$, $c$, work at a competition. Each presses a button if
they accept the attempt. Lamp $y$ must light up if \emph{at least two}
judges accepted the attempt. Such a function is called a
\term{majority function} (from Latin \emph{major}---greater).

\textbf{Steps 1--2. Truth table.} Three inputs give $2^3 = 8$ rows:

\begin{table}[H]
\centering
\caption{Truth table of the majority function}
\label{tab:maj}
\begin{tabular}{c|ccc|c|l}
\toprule
\textbf{No.} & $a$ & $b$ & $c$ & $y$ & \textbf{Term} \\
\midrule
0 & 0 & 0 & 0 & 0 & \\
1 & 0 & 0 & 1 & 0 & \\
2 & 0 & 1 & 0 & 0 & \\
3 & 0 & 1 & 1 & \textbf{1} & $\overline{a}bc$ \\
4 & 1 & 0 & 0 & 0 & \\
5 & 1 & 0 & 1 & \textbf{1} & $a\overline{b}c$ \\
6 & 1 & 1 & 0 & \textbf{1} & $ab\overline{c}$ \\
7 & 1 & 1 & 1 & \textbf{1} & $abc$ \\
\bottomrule
\end{tabular}
\end{table}

\textbf{Step 3. Expression.} For each row where $y=1$, write the product of
all the inputs: an input is taken as is if it equals 1, and inverted if it
equals 0. Such a product equals one \emph{only} in that row. Adding up all
such products gives the function. This form is called the \term{canonical sum
of products} (canonical SOP):
\[
  y = \overline{a}bc + a\overline{b}c + ab\overline{c} + abc.
\]

\textbf{Step 4. Simplification.} The term $abc$ can be repeated as many times
as we like ($x + x = x$). Add it two more times and group:
\[
  y = (\overline{a}bc + abc) + (a\overline{b}c + abc) + (ab\overline{c} + abc)
    = bc(\overline{a}+a) + ac(\overline{b}+b) + ab(\overline{c}+c)
    = bc + ac + ab.
\]
Instead of four three-input AND gates and three inverters, we need only
three two-input AND gates and one OR gate.

\textbf{Step 5. Circuit.}

\begin{figure}[H]
\centering
\begin{tikzpicture}
  \foreach \n/\x in {a/0, b/0.6, c/1.2} {
    \draw[very thick, FpgaBlue] (\x,0.6) node[above, text=FpgaDark] {$\n$} -- (\x,-3.4);
  }
  \node[gAND] (g1) at (3.4,0) {};
  \node[gAND] (g2) at (3.4,-1.4) {};
  \node[gAND] (g3) at (3.4,-2.8) {};
  \node[gOR, logic gate inputs=nnn, minimum height=14mm] (o) at (7.4,-1.4) {}; % fig-ok: 3-input OR 14 mm: inputs not closer than 4.5 mm
  \draw[wire] (g1.input 1) -- (0,0 |- g1.input 1) node[circle, fill, inner sep=1.2pt] {};
  \draw[wire] (g1.input 2) -- (0.6,0 |- g1.input 2) node[circle, fill, inner sep=1.2pt] {};
  \draw[wire] (g2.input 1) -- (0,0 |- g2.input 1) node[circle, fill, inner sep=1.2pt] {};
  \draw[wire] (g2.input 2) -- (1.2,0 |- g2.input 2) node[circle, fill, inner sep=1.2pt] {};
  \draw[wire] (g3.input 1) -- (0.6,0 |- g3.input 1) node[circle, fill, inner sep=1.2pt] {};
  \draw[wire] (g3.input 2) -- (1.2,0 |- g3.input 2) node[circle, fill, inner sep=1.2pt] {};
  \draw[wire] (g1.output) -| ([xshift=-4mm]o.input 1) -- (o.input 1);
  \draw[wire] (g2.output) -- (o.input 2);
  \draw[wire] (g3.output) -| ([xshift=-4mm]o.input 3) -- (o.input 3);
  \node[lbl, anchor=south west] at ([xshift=1mm]g1.output) {$ab$};
  \node[lbl, anchor=south west] at ([xshift=1mm]g2.output) {$ac$};
  \node[lbl, anchor=south west] at ([xshift=1mm]g3.output) {$bc$};
  \draw[wire] (o.output) -- ++(0.8,0) node[right] {$y$};
\end{tikzpicture}
\caption{The ``2 of 3'' majority gate: $y = ab + ac + bc$}
\label{fig:maj}
\end{figure}

\subsection{Example: A Guard with Night Mode}

Now a word problem that will be useful to us later. An alarm must sound if
the door is open \emph{and} (it is night \emph{or} the guard is on). Name the
inputs: $d$---the door is open, $n$---night, $s$---the guard is on. The
parentheses in the condition show the order: first ``night or guard,'' then
``and door'':
\[
  y = d\cdot(n + s).
\]
If we expand the parentheses we get $y = dn + ds$: two AND gates and one OR
against one OR and one AND in the original form. The parenthesized form is
\emph{cheaper} (two gates instead of three) and also more faithful to the
meaning: it repeats the condition word for word. How to choose between such
forms we will examine in detail in the next chapter.

\begin{diybox}[From words to a circuit]
\textbf{Task.} \stars{1} Build the ``2 of 3'' circuit from
Figure~\ref{fig:maj} and check all 8 combinations. \stars{2} Change the
condition: the lamp must light up if \emph{exactly two} accepted. Build the
table and the expression. \stars{3} A stage of the ``Security Panel''
project: add an arming input $arm$ to ``2 of 3'': the alarm sounds only when
$arm=1$ and at least two of the three door sensors agree.

\textbf{In Logra.} \stars{1} \emph{Learn} tab: the lesson
\emph{``Two out of three''} (book ``Logic gates''); the example
\code{majority}. \stars{2} The example \code{exactly-two}: compare it with
your own expression. \stars{3} \emph{Learn} tab, book ``Projects,'' the
lesson \emph{A ``2 of 3'' vote with a guard}.

\textbf{Expected result.} In \stars{1} all 8 rows of Table~\ref{tab:maj}
match. In \stars{2} the lamp is off for $111$. In \stars{3} with $arm=0$ the
lamp is off for any $a,b,c$; the output is one in exactly four rows out of
sixteen ($arm\,abc=1011$, $1101$, $1110$, $1111$).

\textbf{How to check.} \emph{Waveforms} tab: set $a,b,c$ in turn as a
counter 0--7 and compare $y$ with the table; the lesson is checked by the
built-in simulator (the check button). Acceptance criterion for stage~3: four
rows with $y=1$ out of sixteen and the alarm off when $arm=0$. We will find
the minimal form of this function in Chapter~\ref{ch:minimize}.
\end{diybox}

\section{Karnaugh Maps}

Simplifying expressions algebraically is possible, but it is easy to make a
mistake. For functions of 2--4 variables there is a visual method---the
\term{Karnaugh map}. It is the same truth table, but drawn as a rectangle so
that \emph{neighboring cells differ in the value of only one variable}. That
is why the columns are labeled in the order 00, 01, 11, 10 (Gray code) and not
00, 01, 10, 11.

The simplification rule: group the ones into rectangular groups of 2, 4, 8
cells (a power of two), trying to make the groups as large as possible. Each
group corresponds to one product, in which only the variables that
\emph{do not change} within the group remain.

\begin{figure}[H]
\centering
\begin{tikzpicture}[cell/.style={draw, minimum size=11mm, font=\sffamily\large}]
  \foreach \c/\lab in {0/00, 1/01, 2/11, 3/10} \node[font=\sffamily\small] at (\c*1.1,1.05) {\lab};
  \foreach \r/\lab in {0/0, 1/1} \node[font=\sffamily\small] at (-0.95,-\r*1.1) {\lab};
  \node[font=\sffamily\small] at (-1.1,1.05) {$a \backslash bc$};
  % values: row a=0: bc=00,01,11,10 -> 0,0,1,0; a=1: 0,1,1,1
  \foreach \r/\c/\v in {0/0/0, 0/1/0, 0/2/1, 0/3/0, 1/0/0, 1/1/1, 1/2/1, 1/3/1}
    {\ifnum\v=1 \def\cc{FpgaLight}\else\def\cc{white}\fi
     \node[cell, fill=\cc] at (\c*1.1,-\r*1.1) {\v};}
  \draw[very thick, FpgaRed, rounded corners=4pt] (1.7,-1.62) rectangle (2.7,0.52);
  \draw[very thick, FpgaGreen, rounded corners=4pt] (0.6,-1.56) rectangle (2.62,-0.64);
  \draw[very thick, FpgaOrange, rounded corners=4pt] (1.78,-1.48) rectangle (3.8,-0.72);
  \node[font=\sffamily\small, text=FpgaRed, anchor=west] at (4.6,0.2) {$bc$ --- column $bc=11$ ($a$ changes)};
  \node[font=\sffamily\small, text=FpgaGreen, anchor=west] at (4.6,-0.5) {$ac$ --- cells $abc=101, 111$};
  \node[font=\sffamily\small, text=FpgaOrange, anchor=west] at (4.6,-1.2) {$ab$ --- cells $abc=111, 110$};
  \node[font=\sffamily\small\bfseries, anchor=west] at (4.6,-2.0) {$y = ab + ac + bc$};
\end{tikzpicture}
\caption{Karnaugh map of the majority function: three groups of two cells
give three terms}
\label{fig:karnaugh}
\end{figure}

Let us see where the terms come from. The red group is the column $bc=11$: in
its two cells only $a$ changes, while $b=1$ and $c=1$ stay, so the term is
$bc$. The green group is cells $101$ and $111$: here $b$ changes, $a=1$ and
$c=1$ stay, and the term is $ac$. The orange group: $c$ changes, $a=1$ and
$b=1$ stay, and the term is $ab$. The three groups give the same three terms
that we got by algebra, but without the calculations.

Two features of the three-variable map that people forget. First, the map is
\emph{wrapped around horizontally}: the outer columns $00$ and $10$ are also
neighbors (they differ only in $b$). Second, the same one may belong to
several groups: overlapping groups are allowed and often useful, because they
make the groups larger.

\begin{infobox}[Who will simplify circuits in large projects?]
When circuits become large, the simplification is done by a special program,
and much better than a human can do it. But understanding \emph{how} it is
done is still useful: it helps you estimate how many resources your circuit
will take and how fast it will run. The full set of techniques (four
variables, don't-care states, the formal algorithm) is in the next chapter.
\end{infobox}

\begin{diybox}[A Karnaugh map on paper and in a circuit]
\textbf{Task.} \stars{1} For $F_2=\Sigma m(1,3,5,7)$ draw the map, find the
minimal form and build it in the Designer. \stars{2} For
$F_3=\Sigma m(3,5,6,7)$ find the minimal form, build it and compare it with
the ``2 of 3'' circuit (is it the same function written differently?).
\stars{3} Invent a function of three variables whose minimal form consists of
three terms of two letters each, but whose groups do \emph{not} include any
``wrap-around across the edge.''

\textbf{In Logra.} \emph{Designer} tab: a circuit of AND, OR, NOT blocks. Set
the inputs $a,b,c$ with switches.

\textbf{Expected result.} For $F_2$ the minimal form is a single wire,
$F_2=c$; in the circuit there is not a single gate between the input and the
output. For $F_3$ you get $ab+ac+bc$, that is, the same majority function.

\textbf{How to check.} \emph{Waveforms} tab: go through all 8 combinations;
the output must match the column $F$ of your map. \emph{RTL} tab: compare your
circuit with the figure in this chapter.
\end{diybox}

\section{A Truth Table in Verilog}
\label{sec:comb-verilog}

A circuit does not have to be drawn: it can be \emph{written}. The Verilog
language (Chapter~\ref{ch:verilog}) describes the same function as text, and
the synthesizer builds the circuit by itself. Combinational logic needs two
tools: \code{assign} for an expression and \code{always @*} with \code{case}
for a table. Let us take the majority function ``2 of 3''
(Table~\ref{tab:maj}) and write the same thing in three ways.

\textbf{Way 1: from the canonical SOP.} Each table row with a one is one
term.

\vfile{\codepath/b3_comb/maj3_sdnf.v}{code/b3\_comb/maj3\_sdnf.v}

\textbf{How to read it.} The sign \code{\textasciitilde} is NOT, \code{\&} is AND,
\code{|} is OR; line by line it repeats the terms
$\overline{a}bc+a\overline{b}c+ab\overline{c}+abc$. The output \code{y} is updated
on any change of \code{a}, \code{b}, \code{c}: this is not a ``command'' but a
permanent connection, like a wire.

\textbf{Way 2: the minimal form from the Karnaugh map.}

\vfile{\codepath/b3_comb/maj3_min.v}{code/b3\_comb/maj3\_min.v}

The same result with three terms instead of four and no inversions: this is
exactly what simplification buys.

\textbf{Way 3: a table in \code{case}.} The input is concatenated into a
three-bit number \code{\{a, b, c\}}, and for each table row the value of the
output is written.

\vfile{\codepath/b3_comb/maj3_case.v}{code/b3\_comb/maj3\_case.v}

\textbf{How to read it.} The \code{always @*} block is re-evaluated whenever any
input changes; \code{reg} here is not a ``register'', just a variable assigned
inside the block. All 8 rows are listed, so the output is always defined. This
form is the longest, but nothing has to be simplified: the truth table is
rewritten ``as is''.

The three forms describe one function although they look different (how the
minimal form relates to \code{case} we will examine in
Chapter~\ref{ch:minimize}). The testbench \code{tb\_maj3.v} checks this: it
sweeps all 16 combinations and compares each module with the reference
$a+b+c\ge2$.

\vfile{\codepath/b3_comb/tb_maj3.v}{code/b3\_comb/tb\_maj3.v}

\textbf{How to run it in Logra.} The \emph{Code} tab: paste the module
\code{maj3\_sdnf}, press \emph{Synthesize} (Ctrl+Enter). The \emph{RTL} tab
shows the circuit; on the \emph{Waveforms} tab drive $a,b,c$ with a counter
$0$--$7$ and compare $y$ with Table~\ref{tab:maj}. Then replace the body of
the module with the second or the third form and make sure the waveform has not
changed.

\textbf{A guard with night mode} is also written in one line: the brackets
repeat the condition.

\vfile{\codepath/b3_comb/guard_night.v}{code/b3\_comb/guard\_night.v}

\begin{warnbox}[Typical mistakes in Verilog]
\begin{itemize}
  \item \textbf{Rows missing in \code{case} and no \code{default}.} For the
        uncovered input values the output ``remembers'' the past: a latch
        appears instead of a combinational circuit. List all rows or add
        \code{default}.
  \item \textbf{Not assigned in all branches of \code{if}.} The same thing: if
        under some condition the output is not assigned, a latch appears.
        Assign a default value at the beginning of the block.
  \item \textbf{\code{<=} in a combinational block.} For combinational logic
        write \code{=}; keep \code{<=} for flip-flops
        (Chapter~\ref{ch:ff}).
\end{itemize}
\end{warnbox}

\begin{diybox}[Complete and change a module]
\textbf{Task.} \stars{1} \emph{Repeat.} Type \code{maj3\_min} on the
\emph{Code} tab and check all 8 rows on the waveform. \stars{2}
\emph{Change it.} Make the lamp ``exactly two of three'' (for $111$ the lamp
is off), without changing the ports. \stars{3} \emph{Complete it (a stage of
the ``Security Panel'' project).} Add an input $arm$ to ``2 of 3'': the alarm
only when $arm=1$.

\textbf{In Logra.} The \emph{Code} tab, then \emph{Synthesize}; the
\emph{Waveforms} tab: for \stars{3} drive $arm\,a\,b\,c$ with a counter
$0$--$15$.

\textbf{Expected result.} \stars{2} A one exactly in the rows $011,101,110$.
\stars{3} For $arm=0$ the output is zero on all 8 combinations of $abc$; for
$arm=1$ the output matches ``2 of 3'': four ones out of sixteen.

\textbf{How to check.} Compare the waveform with the table column; run
\code{tb\_maj3.v} (it checks the model answers below and prints \code{OK}).
\end{diybox}

\textbf{Model answer to \stars{2}} (count the ones and compare with two):

\vfile{\codepath/b3_comb/exactly2.v}{code/b3\_comb/exactly2.v}

\textbf{Model answer to \stars{3}:}

\vfile{\codepath/b3_comb/guard_arm.v}{code/b3\_comb/guard\_arm.v}

\section{Propagation Delay}

So far we have assumed that a gate responds instantly. In fact transistors
need time to switch, and wires need time to charge. The output of a gate
changes a time $t_\text{pd}$ after the input. This is the
\term{propagation delay}. In modern chips it ranges from fractions of a
nanosecond to a few nanoseconds.

\begin{figure}[H]
\centering
\begin{tikzpicture}[x=1cm, y=1cm]
  \draw[very thick, FpgaBlue] (0,1.6) -- (2,1.6) -- (2.1,2.4) -- (6,2.4) -- (6.1,1.6) -- (9,1.6);
  \draw[very thick, FpgaOrange] (0,0.8) -- (3,0.8) -- (3.1,0) -- (7,0) -- (7.1,0.8) -- (9,0.8);
  \node[font=\sffamily\small, anchor=east] at (-0.1,2.0) {$a$};
  \node[font=\sffamily\small, anchor=east] at (-0.1,0.4) {$\overline{a}$};
  \draw[dashed, FpgaGray] (2.05,2.6) -- (2.05,-0.4); \draw[dashed, FpgaGray] (3.05,2.6) -- (3.05,-0.4);
  \draw[<->, thick] (2.05,-0.3) -- (3.05,-0.3) node[midway, below, font=\sffamily\small] {$t_\text{pd}$};
  \draw[dashed, FpgaGray] (6.05,2.6) -- (6.05,-0.4); \draw[dashed, FpgaGray] (7.05,2.6) -- (7.05,-0.4);
  \draw[<->, thick] (6.05,-0.3) -- (7.05,-0.3) node[midway, below, font=\sffamily\small] {$t_\text{pd}$};
\end{tikzpicture}
\caption{The inverter output changes after a delay $t_\text{pd}$}
\end{figure}

If a signal passes through several gates in a row, the delays add up. The
longest path from an input to an output is called the \term{critical path}.
It determines how fast the circuit can run: until the correct value has
settled at the output, the result must not be read.

\begin{figure}[H]
\centering
\begin{tikzpicture}
  \node[gAND] (g1) at (0,0) {};
  \node[gOR] (g2) at (2.6,-1.1) {};
  \node[gXOR] (g3) at (5.2,-2.2) {};
  \node[gNOT] (g4) at (7.8,-2.2) {};
  \draw[wire, FpgaRed, very thick] (g1.input 1) -- ++(-0.5,0) coordinate (pa) node[left, text=FpgaDark] {$a$};
  \draw[wire] (g1.input 2) -- ++(-0.5,0) node[left] {$b$};
  \draw[wire, FpgaRed, very thick] (g1.output) -| ([xshift=-4mm]g2.input 1) -- (g2.input 1);
  \draw[wire] (g2.input 2) -- (g2.input 2 -| pa) node[left] {$c$};
  \draw[wire, FpgaRed, very thick] (g2.output) -| ([xshift=-4mm]g3.input 1) -- (g3.input 1);
  \draw[wire] (g3.input 2) -- (g3.input 2 -| pa) node[left] {$d$};
  \draw[wire, FpgaRed, very thick] (g3.output) -- (g4.input);
  \draw[wire, FpgaRed, very thick] (g4.output) -- ++(0.6,0) node[right, text=FpgaDark] {$y$};
  \node[font=\sffamily\small, text=FpgaRed, anchor=west] at (10.2,-1.9) {critical path:};
  \node[font=\sffamily\small, text=FpgaRed, anchor=west] at (10.2,-2.5) {$t = 4\,t_\text{pd}$};
\end{tikzpicture}
\caption{The critical path passes through four gates}
\end{figure}

A simple consequence: the benefit of simplifying a circuit lies not only in
the number of gates but also in time. If a chain of four gates can be
replaced by two, the critical path is halved.

\subsection{Races and Glitches}

Because of the delays, short false pulses---\term{glitches}---can appear at
the output of a combinational circuit. A classic example:
$y = a\cdot\overline{a}$. By the laws of algebra $y$ is always 0. But the
inverter responds with a delay, and for a short time after an edge of $a$
\emph{both} inputs of the AND gate are 1.

\begin{figure}[H]
\centering
\begin{tikzpicture}
  \node[gAND] (g) at (4.0,0) {};
  \node[gNOT] (n) at (1.8,-1.2) {};
  \draw[wire] (0,0 |- g.input 1) node[left] {$a$} -- (g.input 1);
  \fill (0.5,0 |- g.input 1) circle (1.3pt);
  \draw[wire] (0.5,0 |- g.input 1) |- (n.input);
  \draw[wire] (n.output) -| ([xshift=-4mm]g.input 2) -- (g.input 2);
  \draw[wire] (g.output) -- ++(0.5,0) node[right] {$y$};
  \begin{scope}[xshift=7.4cm, yshift=0.9cm, x=1cm, y=1cm]
    \draw[very thick, FpgaBlue] (0,0) -- (1.5,0) -- (1.55,0.5) -- (5,0.5);
    \draw[very thick, FpgaOrange] (0,-0.4) -- (2.4,-0.4) -- (2.45,-0.9) -- (5,-0.9);
    \draw[very thick, FpgaRed] (0,-1.8) -- (1.6,-1.8) -- (1.65,-1.3) -- (2.5,-1.3) -- (2.55,-1.8) -- (5,-1.8);
    \node[font=\sffamily\small, anchor=east] at (-0.1,0.25) {$a$};
    \node[font=\sffamily\small, anchor=east] at (-0.1,-0.65) {$\overline{a}$};
    \node[font=\sffamily\small, anchor=east] at (-0.1,-1.55) {$y$};
    \node[font=\sffamily\small, text=FpgaRed, anchor=west] at (2.7,-1.4) {false pulse!};
  \end{scope}
\end{tikzpicture}
\caption{How a glitch arises from the inverter's delay}
\end{figure}

This is an example of a \emph{race}: two paths from the same input to the
output (the direct one and the one through the inverter) have different
lengths, and the signals ``run a race'' along them. For a moment the output
takes a value that is not in the truth table.

\subsubsection{Glitches and the Karnaugh map}

A glitch can appear even in a circuit whose logic is flawless. Take
$F = AB + \overline{A}C$ and fix $B=C=1$. Then $F = A + \overline{A}$ must
always equal 1. But the term $AB$ turns off when $A$ goes from 1 to 0, while
the term $\overline{A}C$ turns on later, because the signal $\overline{A}$ is
still ``on its way'' through the inverter. For a short time \emph{both} terms
are off, and a zero glitch appears at the output $F$.

On the Karnaugh map this is visible at once: the two groups $AB$ and
$\overline{A}C$ \emph{touch} but do not overlap. The transition between them
goes along the boundary where neither group holds the one. The cure is to add
an overlapping group $BC$: it covers both cells at the boundary and holds the
output at one while $A$ switches.

\begin{figure}[H]
\centering
\begin{tikzpicture}[cell/.style={draw, minimum size=11mm, font=\sffamily\large}]
  \foreach \c/\lab in {0/00, 1/01, 2/11, 3/10} \node[font=\sffamily\small] at (\c*1.1,1.05) {\lab};
  \foreach \r/\lab in {0/0, 1/1} \node[font=\sffamily\small] at (-0.95,-\r*1.1) {\lab};
  \node[font=\sffamily\small] at (-1.1,1.05) {$A \backslash BC$};
  \foreach \r/\c/\v in {0/0/0, 0/1/1, 0/2/1, 0/3/0, 1/0/0, 1/1/0, 1/2/1, 1/3/1}
    {\ifnum\v=1 \def\cc{FpgaLight}\else\def\cc{white}\fi
     \node[cell, fill=\cc] at (\c*1.1,-\r*1.1) {\v};}
  \draw[very thick, FpgaOrange, rounded corners=4pt] (0.65,0.42) rectangle (2.55,-0.42);
  \draw[very thick, FpgaGreen, rounded corners=4pt] (1.85,-0.68) rectangle (3.75,-1.52);
  \draw[very thick, dashed, FpgaRed, rounded corners=4pt] (1.9,0.5) rectangle (2.5,-1.6);
  \node[font=\sffamily\small, text=FpgaOrange, anchor=west] at (4.6,0.2) {$\overline{A}C$ --- cells $A=0$};
  \node[font=\sffamily\small, text=FpgaGreen, anchor=west] at (4.6,-0.5) {$AB$ --- cells $A=1$};
  \node[font=\sffamily\small, text=FpgaRed, anchor=west] at (4.6,-1.2) {$BC$ --- added against the glitch};
  \node[font=\sffamily\small\bfseries, anchor=west] at (4.6,-2.0) {$F = AB + \overline{A}C + BC$};
\end{tikzpicture}
\caption{The groups $\overline{A}C$ and $AB$ touch but do not overlap.
The dashed group $BC$ removes the glitch, although the logic does not need it}
\label{fig:hazard_karnaugh}
\end{figure}

A paradox: the added group is \emph{redundant} from the point of view of the
minimal form, yet it is needed for reliable operation. So the ``minimal'' and
the ``best'' circuits can differ, and this is worth remembering.

\begin{warnbox}[How to deal with this]
Glitches are the main reason why, in serious circuits, the outputs of
combinational logic are ``read'' not at random moments but only on the edge
of the clock signal, when all transient processes are over. This is called
\term{synchronous design} (Chapter~\ref{ch:sync}). A practical rule: a glitch
is harmless if nobody reads it until it is over.
\end{warnbox}

\begin{diybox}[Delay and glitch]
\textbf{Task.} \stars{1} For the circuit from the ``Critical path'' figure,
calculate the delay if $t_\text{pd}=1.2$~ns for each gate, and draw the
timing diagram of $y=a\cdot\overline a$ with the inverter delay, marking the
glitch. \stars{2} For $F=AB+\overline{A}C$ with $B=C=1$ draw the diagram and
find the moment when $F$ briefly drops to 0. Add the term $BC$ and repeat.
\stars{3} Invent a circuit in which a glitch is harmful (hint: the output is
connected to an input that reacts to any short pulse).

\textbf{In Logra.} \emph{Designer} tab: build $y=a\cdot\overline a$ (input
$a$, a NOT block, an AND block), and next to it a chain of five NOT blocks.
Then the example \code{button-debounce} (the example library): how contact
bounce resembles a glitch.

\textbf{Expected result.} For the calculation: $4\cdot 1.2 = 4.8$~ns. Five
inversions give $y=\overline a$. On the diagram the output $y=a\cdot\overline
a$ stays zero for any $a$: ordinary logic simulation does not model gate
delays. This is not a bug but an important observation: the glitch is not
visible on such a diagram, yet it exists in a real circuit.

\textbf{How to check.} \emph{Waveforms} tab: make sure that $y\equiv0$ (the
logic is correct) and compare with your paper diagram with the glitch. The
difference between them is exactly the problem for which circuits with memory
use a common clock signal. Where the nanoseconds come from is covered in
Chapter~\ref{ch:electrical}.
\end{diybox}

\begin{warnbox}[Common mistakes]
\begin{itemize}
  \item \textbf{Simplifying ``however it comes out.''} Algebraic
        transformations are easy to spoil: you skip an inversion or forget
        that $\overline{a+b}\neq\overline a+\overline b$. After each
        simplification check the result against the truth table on at least
        three rows.
  \item \textbf{Grouping the wrong cells.} A group must be a rectangle of
        $1,2,4,8$ cells. Three cells in a row, an ``L-shape,'' and a diagonal
        are not groups.
  \item \textbf{Forgetting that the edges are neighbors.} The columns $00$ and
        $10$ of the map are neighbors. If you found two ones at the edges and
        did not combine them, you got an extra term.
  \item \textbf{Treating a glitch as a logic error.} A glitch does not
        violate the truth table: it lives in the transition moment and is
        removed by clocking or by adding an overlapping group.
\end{itemize}
\end{warnbox}

\begin{ideabox}
A combinational circuit has no memory. It is designed from a truth table and
simplified algebraically or with a Karnaugh map. Real gates have a delay,
which makes short false pulses possible.
\end{ideabox}

\section{Questions and Problems}
\label{sec:tasks-comb}

Answers are collected in Appendix~\ref{app:answers}.

\begin{taskbox}[Chapter problems]
\begin{enumerate}
  \item \stars{1}~Find the minimal form of the functions $F_1=\Sigma m(0,1,4,5)$, $F_2=\Sigma m(1,3,5,7)$ and $F_3=\Sigma m(3,5,6,7)$ of three variables $A,B,C$ ($A$ is the most significant bit) using a Karnaugh map.
  \item \stars{1}~A signal passes through 5 gates with a delay of 1.5~ns each. What is the delay of this path? How is it related to the maximum frequency of the circuit?
  \item \stars{2}~An alarm sounds if the door is open \emph{and} (it is night \emph{or} the guard is on). Write the function and draw the circuit.
  \item \stars{1}~Which of these blocks are combinational and which are sequential: an adder, a counter, a multiplexer, a button that toggles a lamp on every press?
  \item \stars{2}~For the function $F = AB + \overline{A}C$ with $B=C=1$, $A$ changes. Does a glitch occur? Which term must be added to remove it?
  \item \stars{2}~\emph{Find the error.} A student simplified $y=\overline{a}bc+a\overline{b}c+ab\overline{c}+abc$ and got $y=ab+c$. Check the answer against the truth table and find the row where the student went wrong.
  \item \stars{2}~\emph{Find the error.} On the three-variable map for $F=\Sigma m(1,2,3)$ a student circled an ``L-shape'' of three cells $1,3,2$ as one group and wrote $F=\overline{A}$. What is the student's function in reality? What should the minimal form be?
  \item \stars{3}~\emph{A project with a criterion.} Build the ``guard with night mode'' circuit $y=d\,(n+s)$ in the Designer and make sure that for all 8 combinations $y$ matches the table. Acceptance criterion: 8 out of 8 rows match on the Waveforms tab.
  \item \stars{3}~\emph{A project with a criterion (a stage of the ``Security Panel'' project).} Add an arming input $arm$ to the ``2 of 3'' circuit and build the minimal form $y=arm\cdot(ab+ac+bc)$. Acceptance criterion: with $arm=0$ the output is zero for any $a,b,c$; with $arm=1$ the output matches Table~\ref{tab:maj}. How many gates did you get?
  \item \stars{3}~Show that for the function $y=ab+\overline{a}c$ a glitch occurs at $b=c=1$, while for the function $y=ab+\overline{a}c+bc$ it does not. How is this related to the Karnaugh map?
  \item \stars{2}~\emph{Verilog.} Write the function ``exactly two of three'' with a single \code{assign} and check it on all 8 combinations on the Waveforms tab.
  \item \stars{2}~\emph{Find the error (Verilog).} A student described ``2 of 3'' with \code{case (\{a,b,c\})} but listed only the four rows with a one and wrote no \code{default}. What happens on the other four rows, and how to fix it?
\end{enumerate}
\end{taskbox}

\section{What's Next}

We can build a circuit from a table and simplify a function of three
variables on a Karnaugh map. But there are far more functions of four
variables, some combinations have no value at all, and real programs simplify
functions of dozens of variables. In the next chapter we will study the
four-variable Karnaugh map, don't-care states, incorrect groupings, and see
how a minimal form turns into a real circuit in Logra.
